
\documentclass[a4paper,11pt]{article}
\usepackage[a4paper, margin=8em]{geometry}

% usa i pacchetti per la scrittura in italiano
\usepackage[french,italian]{babel}
\usepackage[T1]{fontenc}
\usepackage[utf8]{inputenc}
\frenchspacing 

% usa i pacchetti per la formattazione matematica
\usepackage{amsmath, amssymb, amsthm, amsfonts}

% usa altri pacchetti
\usepackage{gensymb}
\usepackage{hyperref}
\usepackage{standalone}

% imposta il titolo
\title{Appunti Comunicazioni Numeriche}
\author{Luca Seggiani}
\date{2026}

% imposta lo stile
% usa helvetica
\usepackage[scaled]{helvet}
% usa palatino
\usepackage{palatino}
% usa un font monospazio guardabile
\usepackage{lmodern}

\renewcommand{\rmdefault}{ppl}
\renewcommand{\sfdefault}{phv}
\renewcommand{\ttdefault}{lmtt}

% disponi il titolo
\makeatletter
\renewcommand{\maketitle} {
	\begin{center} 
		\begin{minipage}[t]{.8\textwidth}
			\textsf{\huge\bfseries \@title} 
		\end{minipage}%
		\begin{minipage}[t]{.2\textwidth}
			\raggedleft \vspace{-1.65em}
			\textsf{\small \@author} \vfill
			\textsf{\small \@date}
		\end{minipage}
		\par
	\end{center}

	\thispagestyle{empty}
	\pagestyle{fancy}
}
\makeatother

% disponi teoremi
\usepackage{tcolorbox}
\newtcolorbox[auto counter, number within=section]{theorem}[2][]{%
	colback=blue!10, 
	colframe=blue!40!black, 
	sharp corners=northwest,
	fonttitle=\sffamily\bfseries, 
	title=~\thetcbcounter: #2, 
	#1
}

% disponi definizioni
\newtcolorbox[auto counter, number within=section]{definition}[2][]{%
	colback=red!10,
	colframe=red!40!black,
	sharp corners=northwest,
	fonttitle=\sffamily\bfseries,
	title=~\thetcbcounter: #2,
	#1
}

% U.D.M
\newcommand{\amp}{\ensuremath{\mathrm{A}}}
\newcommand{\volt}{\ensuremath{\mathrm{V}}}
\newcommand{\meter}{\ensuremath{\mathrm{m}}}
\newcommand{\second}{\ensuremath{\mathrm{s}}}
\newcommand{\farad}{\ensuremath{\mathrm{F}}}
\newcommand{\henry}{\ensuremath{\mathrm{H}}}
\newcommand{\siemens}{\ensuremath{\mathrm{S}}}

% circuiti
\usepackage{circuitikz}
\usetikzlibrary{babel}

% disegni
\usepackage{pgfplots}
\pgfplotsset{width=10cm,compat=1.9}

% disponi codice
\usepackage{listings}
\usepackage[table]{xcolor}

\lstdefinestyle{codestyle}{
		backgroundcolor=\color{black!5}, 
		commentstyle=\color{codegreen},
		keywordstyle=\bfseries\color{magenta},
		numberstyle=\sffamily\tiny\color{black!60},
		stringstyle=\color{green!50!black},
		basicstyle=\ttfamily\footnotesize,
		breakatwhitespace=false,         
		breaklines=true,                 
		captionpos=b,                    
		keepspaces=true,                 
		numbers=left,                    
		numbersep=5pt,                  
		showspaces=false,                
		showstringspaces=false,
		showtabs=false,                  
		tabsize=2
}

\lstdefinestyle{shellstyle}{
		backgroundcolor=\color{black!5}, 
		basicstyle=\ttfamily\footnotesize\color{black}, 
		commentstyle=\color{black}, 
		keywordstyle=\color{black},
		numberstyle=\color{black!5},
		stringstyle=\color{black}, 
		showspaces=false,
		showstringspaces=false, 
		showtabs=false, 
		tabsize=2, 
		numbers=none, 
		breaklines=true
}

\lstdefinelanguage{javascript}{
	keywords={typeof, new, true, false, catch, function, return, null, catch, switch, var, if, in, while, do, else, case, break},
	keywordstyle=\color{blue}\bfseries,
	ndkeywords={class, export, boolean, throw, implements, import, this},
	ndkeywordstyle=\color{darkgray}\bfseries,
	identifierstyle=\color{black},
	sensitive=false,
	comment=[l]{//},
	morecomment=[s]{/*}{*/},
	commentstyle=\color{purple}\ttfamily,
	stringstyle=\color{red}\ttfamily,
	morestring=[b]',
	morestring=[b]"
}

% disponi sezioni
\usepackage{titlesec}

\titleformat{\section}
	{\sffamily\Large\bfseries} 
	{\thesection}{1em}{} 
\titleformat{\subsection}
	{\sffamily\large\bfseries}   
	{\thesubsection}{1em}{} 
\titleformat{\subsubsection}
	{\sffamily\normalsize\bfseries} 
	{\thesubsubsection}{1em}{}

% disponi alberi
\usepackage{forest}

\forestset{
	rectstyle/.style={
		for tree={rectangle,draw,font=\large\sffamily}
	},
	roundstyle/.style={
		for tree={circle,draw,font=\large}
	}
}

% disponi algoritmi
\usepackage{algorithm}
\usepackage{algorithmic}
\makeatletter
\renewcommand{\ALG@name}{Algoritmo}
\makeatother

% disponi numeri di pagina
\usepackage{fancyhdr}
\fancyhf{} 
\fancyfoot[L]{\sffamily{\thepage}}

\makeatletter
\fancyhead[L]{\raisebox{1ex}[0pt][0pt]{\sffamily{\@title \ \@date}}} 
\fancyhead[R]{\raisebox{1ex}[0pt][0pt]{\sffamily{\@author}}}
\makeatother

\begin{document}
% sezione (data)
\section{Lezione del 08-05-26}

% stili pagina
\thispagestyle{empty}
\pagestyle{fancy}

% testo
Alla fine della scorsa lezione abbiamo visto la decodifica a \textit{massima verosomiglianza} nei sistemi di comunicazione. Avevamo quindi dimostrato che questa può rilevare fino a:
$$
d_{min} - 1
$$
errori.

\subsubsection{Correzione degli errori}
Non possiamo mai risolvere tutti gli errori che rileviamo, e anzi quando più della metà degli errori che possiamo rilevare al massimo si verifica, non sappiamo più come correggere la parola di codice. Questo significherà che un codice lineare a blocco può correggere fino a:
$$
\frac{d_{min - 1}}{2}
$$
errori.

\subsubsection{Scrambling}
Notiamo che una prerogativa che avevamo nell'analisi matematica della massima verosomiglianza era che la probabilità di errore sul bit fosse distribuita ugualmente, cioè i rumori sui bit fossero incorrelati. Questo non è sempre il caso nei sistemi di comunicazione, in quanto spesso gli errori si presentano in \textit{burst} localmente correlati.

Una soluzione che si può adottare è quella dello \textbf{scrambling} (o \textit{interleaving}) dei bit trasmessi, ammesso che il ricevitore abbia modo di recuperare l'ordine originale. In questo modo il rumore di un burst viene distribuito su un lasso temporale più grande, e reso effettivamente incorrelato da bit a bit.

\subsubsection{Decodifica a sindrome}
Chiamiamo \textbf{sindrome} di una parola di codice $\mathbf{y}$ il vettore $\mathbf{s}$ ottenuto come:
$$
\mathbf{s} = \mathbf{y} H^T = (\mathbf{x} + \mathbf{e})H^T = \mathbf{x} H^T + \mathbf{e} H^T = \mathbf{e} H^T
$$

La sindrome è composta da $n-k$ cifre binarie, e ciascuna sindrome è associata a $2^k$ pattern di errore ottenuti sommando al vettore $\mathbf{e}$ le $2^k$ parole di codice.

La decodifica a sindrome si compie quindi nei seguenti passaggi:
\begin{enumerate}
\item Si calcola la sindrome:
$$
\mathbf{s} = \mathbf{y} H^T
$$
\item Si associa la sindrome all'errore di peso minimo (detto \textit{coset leader}) corrispondente, cioè l'errore di peso minimo a cui corrisponde la sindrome associata ad $\mathbf{y}$:
$$
\mathbf{s} \rightarrow \mathbf{e}_{CL} (\mathbf{s})
$$
\item Si corregge l'errore sommando l'errore di peso minimo alla $n$-upla $\mathbf{y}$:
$$
\hat{\mathbf{x}} = \mathbf{y} + \mathbf{e}_{CL}(\mathbf{s})
$$
\end{enumerate}

Si può dimostrare che la decodifica a sindrome è equivalente a quella a massima verosomiglianza.

\subsubsection{Prestazioni dei codici a blocco}
Abbiamo detto che la parola $\mathbf{y}$, data da:
$$
\mathbf{y} = \mathbf{x} + \mathbf{e}
$$
risulta errata solo se il canale introduce più di $t$ errori, dove:
$$
t = \frac{d_{min} - 1}{2} \implies d_{min} = 2 t + 1
$$

Ne segue che la probabilità di errore di una parola è:
$$
P_W(e) = P\{ w(e) > t \} = \sum_{j = t + 1}^n \binom{n}{j} p^j (1 - p)^{n - j}
$$
Su questa facciamo le classiche approssimazioni di:
\begin{itemize}
	\item Prendere solo il primo caso, quello a $t + 1$, visto che è il più pobabile;
	\item Trascurare $(1 - p)$ in quanto è pressappoco $1$.
\end{itemize}
Da cui si ottiene:
$$
P_W(e) \approx \binom{n}{t + 1} p^{t + 1}
$$

La probabilità di errore sul bit si approssima quindi come:
$$
P_B(e) \approx \frac{d_{min}}{n} P_W(e)
$$

\end{document}
